\documentclass{article}
\usepackage{graphicx} % Required for inserting images
\usepackage{amsmath}
\usepackage{amsthm}
\usepackage{latexsym} % Required for \Join
\usepackage{booktabs}
\newtheorem{theorem}{Theorem}
\newtheorem{lemma}{Lemma}
\newtheorem{proposition}{Proposition}
\newtheorem{corollary}{Corollary}
\newtheorem{definition}{Definition}
\newtheorem{assumption}{Assumption}
\newtheorem{remark}{Remark}


\title{TempPasteOutput}
\author{Ruben Eschauzier}
\date{September 2026}

\begin{document}

\maketitle
% Lets look at it closely. Consider Pod A has CR with three triple patterns in linear shape:
% $?s <p1> ?o1 $, $?o1 <p2> ?o2$, $?o2 <p3> ?o3$.
% Then given the authoritativeness constraint, the results of this join form an exclusive group iff:
% $?s \in CR, ?o1\in CR, ?o2 \in CR$

% Here $\in$ is used to indicate authoritativeness of CR over the term binding to the variable, as its in its domain and term starts with the domain specification.

% Now we have query execution

Lets look at it closely. Consider Pod A has CR with three triple patterns in linear shape:
$tp1 = ?s <p1> ?o1 $, $tp2 = ?o1 <p2> ?o2$.
Then given the authoritativeness constraint, the results of this join form an exclusive group iff:
$?s \in CR, ?o1\in CR$.

Here $\in$ is used to indicate authoritativeness of CR over the term binding to the variable, as its in its domain and term starts with the domain specification.
\\ \\
Now we have query execution:
$tp1 \Join tp2$, which now admits four subplans, where (A) indicates the authoritative subset of the plan claimed by the CR of Pod A:

$(tp1(A) \Join tp2(A)) \cup ((tp1 \setminus tp1(A)) \Join tp2) \cup ((tp2 \setminus tp2(A)) \Join tp1) \cup (tp1\Join tp2$.

Obviously, the first term is handled by the CR, while the rest is handled by the default operators of the query.

Thus, the base operators should only filter solution mappings iff:
$\mu(tp1) \in (tp1(A) \Join tp2(A))$ or viceversa for tp2.

However, the claim is you can't know if this is true unless you know the full result. However, from construction of the delegated sub-query to CR we know its a strict linear query, no branches or other things. In this case if a CR is authoritative over an inner triple pattern (not the end one), then both subject and object of the triple pattern $\in CR$.
In that case the subject of the next also $\in CR$, so this will always need to be filtered


\subsection*{Authoritativeness of Delegated Joins}

Let Pod $A$ expose a composite resource $CR$ whose query is the path
\[
tp_1 = {?s}\ \langle p_1\rangle\ {?o_1}, \qquad
tp_2 = {?o_1}\ \langle p_2\rangle\ {?o_2}.
\]

A composite resource does not answer for the whole web, nor necessarily for the whole
pod that publishes it. What it answers for is a parameter of the resource, and the entire
development below is relative to it.

\begin{definition}[Authority scope]
A composite resource $CR$ declares an \emph{authority scope} $S$, a set of RDF terms. We
write $t \in CR$ for $t \in S$.
\end{definition}

\begin{assumption}[Scope completeness]
\label{as:scope}
Every triple whose subject lies in $S$ belongs to the dataset over which $CR$ evaluates
its query.
\end{assumption}

Nothing in what follows depends on how $S$ is expressed. For the resources we consider it
is the namespace of the pod publishing the resource, and Assumption~\ref{as:scope} then
coincides with the authoritativeness assumption already made of pods: a pod is the sole
source of triples about the subjects it hosts. Stating the scope separately makes visible
what that assumption fixes, namely a space, and lets the same development cover resources
that claim less than a whole pod.

A resource aggregating a strict sub-space of its pod satisfies
Assumption~\ref{as:scope} only for that sub-space. Delegating to such a resource is still
sound --- it answers what it answers --- but testing mappings against the wider space is
not, since a mapping may be authoritative to the pod and yet outside what the resource
aggregated. An engine that cannot establish the scope of a resource must therefore
decline to prune rather than approximate it.

Since a pod is authoritative exactly for the subjects it hosts, the \emph{anchor} of a
triple pattern is its subject, and the anchors of a set of patterns $B$ are
$\mathit{anchors}(B) = \{\, \mathit{subj}(tp) \mid tp \in B \,\}$. Pod $A$ can evaluate
$B$ completely if and only if every anchor lies in $S$. For the path above,
\[
{?s} \in CR \quad\text{and}\quad {?o_1} \in CR ,
\]
with no condition on $?o_2$: it is the object of the last pattern and therefore not an
anchor. Where these conditions hold, $A$ is the only possible source of the join, so
the delegated evaluation is exclusive --- exclusive per solution mapping rather than per
pattern, since other pods answer the same patterns for other bindings.

\subsection*{Deciding Membership}

Whether $t \in S$ holds is decided differently for the two kinds of term that can occupy
an anchor position.

For an IRI, the scope is a namespace and the test is syntactic: $t$ lies in $S$ when it
falls under the namespace, compared on path segment boundaries so that a namespace does
not capture a sibling that merely shares a prefix of its name. This costs nothing beyond
a string comparison and requires no knowledge of where the triple was read.

A blank node carries no namespace, so nothing about the term itself places it in a scope.
What does place it is provenance: blank nodes are scoped to the document that introduces
them, and no triple outside that document can refer to one. Hence if the document a blank
node came from is one the resource aggregates, both the triples binding it and every
triple that could extend them are in the resource's dataset, and the anchor condition
holds. Membership for blank nodes is therefore decided by source attribution and cannot
be decided any other way --- the syntactic test is not merely unavailable, it is
meaningless. An engine that does not track which document a mapping was read from must
decline to prune whenever an anchor is a blank node.

\begin{remark}[Blank nodes shared beyond the block]
\label{rem:bnode}
Blank node identity is local to the document or result that carries it, so the blank node
a resource returns in its answer is not the same term as the one a client reads from the
underlying document. A blank-node anchor may therefore be pruned only when the anchor
variable does not also occur outside $B$: if it does, the mappings the resource supplies
in place of the pruned ones will not join with the rest of the query. Whether this applies
is decided once, from the query, and not per mapping.
\end{remark}

\subsection*{Decomposition}

Write $tp_i^{A}$ for the restriction of $tp_i$ to mappings satisfying its anchor
condition, and $tp_i^{\lnot A}$ for its complement. Both are decidable locally, as each
pattern binds its own anchor. The join then splits into four disjoint parts, one per
combination of the two conditions:
\[
tp_1 \bowtie tp_2 \;=\;
\underbrace{(tp_1^{A} \bowtie tp_2^{A})}_{\text{delegated to } CR}
\;\uplus\; (tp_1^{A} \bowtie tp_2^{\lnot A})
\;\uplus\; (tp_1^{\lnot A} \bowtie tp_2^{A})
\;\uplus\; (tp_1^{\lnot A} \bowtie tp_2^{\lnot A}).
\]
The first part is answered by $CR$; the remaining three are answered by the engine's
default operators. This is an identity rather than a plan: the engine evaluates the
ordinary join and removes the delegated part, rather than evaluating three additional
joins. For a path of $n$ patterns the decomposition has $2^n$ parts, of which exactly
one is delegated.

\subsection*{Detecting the Delegated Part}

Deciding $\mu \in (tp_1^{A} \bowtie tp_2^{A})$ appears to require the delegated join
itself, but it does not.

\begin{lemma}[Membership]
Let block $B$ be delegated to $A$. Then
$\mu \in B(A) \iff \mu(t) \in CR$ for every $t \in \mathit{anchors}(B)$.
\end{lemma}

Membership is thus a conjunction of independent authority tests, decidable on any
mapping that binds all anchors, with no materialisation of the delegated result. The
lemma does not depend on the block being a path; it holds for any conjunctive block, and
the anchor set is what varies between shapes. It does depend on
Assumption~\ref{as:scope}: the forward direction is exclusivity, which the anchors give,
and the converse is completeness of $CR$ over $S$, which the assumption gives.

\begin{lemma}[Pruning]
A default operator for $tp \in B$ may discard the mappings satisfying the anchor
conditions of $B$ if and only if $\mathit{anchors}(B) \subseteq \mathit{vars}(tp)$.
\end{lemma}

The condition is necessary because a pattern that does not bind every anchor cannot
distinguish the delegated part from the parts that share its tuples. In the path above,
a $tp_1$ mapping binds both $?s$ and $?o_1$ and so decides both conditions: if it
satisfies them it can occur only in the delegated part, and may be pruned. A $tp_2$
mapping binds $?o_1$ but not $?s$; a $tp_2$ tuple with $?o_1 \in CR$ occurs both in the
delegated part and in $tp_1^{\lnot A} \bowtie tp_2^{A}$, the case in which the path
enters $A$ from outside. Pruning it would lose that part.

\begin{center}
\begin{tabular}{lll}
\toprule
block & anchors & prunable patterns \\
\midrule
star at $?x$        & $\{?x\}$                       & all \\
path of length $2$  & $\{?s, ?o_1\}$                 & $tp_1$ only \\
path of length $n \geq 3$ & $\{?s, ?o_1, \dots, ?o_{n-1}\}$ & none \\
\bottomrule
\end{tabular}
\end{center}

A star is the shape for which pruning is always available: its anchor set is a
singleton, and that anchor occurs in every one of its patterns. For path with $N=2$, we can also delegate. Longer paths distribute
their anchors across patterns, so the delegated part remains detectable but the inputs
to the default operators cannot be reduced.

Throughout, soundness rests on Assumption~\ref{as:scope} together with the requirement
that the resource's query is exactly the delegated block: the decomposition prevents
double counting, but only completeness over $S$ licenses removing the delegated part from
the default plan.


\subsection*{Operators, Coverage, and the Uniform Filter}

The decomposition above is stated over patterns, but an engine prunes at operators, and
an operator need not correspond to a single pattern. Let an operator $o$ of the default
plan \emph{cover} a set $\mathit{cov}(o) \subseteq B$ of the block's patterns and emit
mappings over $\mathit{vars}(\mathit{cov}(o))$. Its \emph{depth} is
$|\mathit{cov}(o)|$: an operator scanning a single triple pattern has depth $1$, and an
operator holding the connected sub-block it has joined so far has the depth of that
sub-block. Depth counts patterns covered, not position along the path.

Anchors need not be variables. If some constant anchor is not in $CR$ the block is not
delegable at all and no filter is registered; otherwise every constant anchor is
discharged once, at registration time, and only the \emph{residual} anchors
\[
\mathit{anchors}_V(B) \;=\; \mathit{anchors}(B) \cap \mathcal{V}
\]
remain to be tested per mapping.

\begin{lemma}[Generalised Pruning]
Let $B$ be delegated to $A$ and let $o$ be an operator of the default plan with
$\mathit{cov}(o) \subseteq B$. Discarding from $o$ every mapping that satisfies all
anchor conditions of $B$ is sound for every database if and only if
\[
\mathit{anchors}_V(B) \;\subseteq\; \mathit{vars}(\mathit{cov}(o)).
\]
\end{lemma}

\begin{proof}
($\Leftarrow$) The condition makes every anchor of $B$ bound in any mapping $\mu$ emitted
by $o$, so the anchor conditions are decidable on $\mu$. Let $\mu$ satisfy them and let
$\nu$ be any mapping over $B$ with $\nu \supseteq \mu$. The anchor conditions constrain
only terms bound by $\mu$, so $\nu$ satisfies them as well, and $\nu \in B(A)$ by
Membership. Every solution of the whole query factors through a mapping over $B$, so
every solution reachable through $\mu$ is reachable through some $\nu \in B(A)$, and
those are exactly what $CR$ supplies. Discarding $\mu$ therefore removes no solution
from the union.

($\Rightarrow$) Let $t \in \mathit{anchors}_V(B) \setminus
\mathit{vars}(\mathit{cov}(o))$, so $o$ can test only the proper subset of the anchor
conditions its mappings bind. Choose a database in which some $\nu \supseteq \mu$ has
$\nu(t) \notin CR$. Then $\nu \notin B(A)$, so $CR$ does not return it, while $\mu$ has
been discarded, so the default plan cannot produce it either. The mapping $\nu$ is lost.
\end{proof}

\begin{corollary}[The first pattern is not privileged]
For a path of length $n$, the depth-$1$ operator of $tp_1$ satisfies the condition if
and only if $n \leq 2$, and no other depth-$1$ operator of the path satisfies it for any
$n$.
\end{corollary}

\begin{proof}
Assume every anchor is a variable, so $\mathit{anchors}_V(B) = \{?s, ?o_1, \dots,
?o_{n-1}\}$, of size $n$. Since $\mathit{vars}(tp_1) = \{?s, ?o_1\}$, the inclusion
holds exactly when $n \leq 2$. For $i \geq 2$ we have $?s \notin \mathit{vars}(tp_i) =
\{?o_{i-1}, ?o_i\}$, so the inclusion fails for every $n$.
\end{proof}

At $n = 2$ the single prunable operator is thus the one scanning $tp_1$, the pattern that
enters the domain. The operator scanning $tp_2$ is not prunable, even though $tp_2$ is
the pattern whose anchor condition is the last one to be decided: its mappings leave
$?s$ free, so a $tp_2$ tuple with $?o_1 \in CR$ serves both the delegated part and the
part in which a triple from outside $A$ connects into $?o_1$. Prunability follows the
variables an operator binds, not the direction in which the path is read.

Concretely, at $n = 3$ a $tp_1$ mapping with $?s \in CR$ and $?o_1 \in CR$ still occurs
in $tp_1^{A} \bowtie tp_2^{A} \bowtie tp_3^{\lnot A}$: the path runs two hops inside $A$
and then leaves it, because $A$ is authoritative for the subject $?o_1$ but says nothing
about where its object $?o_2$ lives. That part is not delegated, $CR$ never emits it, and
pruning $tp_1$ loses it.

The corollary is about coverage, not about position in the path. If the entry point is a
constant $\langle a \rangle \in CR$, then $\mathit{anchors}_V(B) = \{?o_1, ?o_2\}$ at
$n = 3$, and the prunable depth-$1$ operator is the one scanning $tp_2$, whose variables
are exactly those two, while $tp_1$, binding only $?o_1$, is not prunable. Binding the
entry point buys back one level of pruning depth, and moves the prunable pattern along
the path.

\subsection*{Methodology}

The lemma is awkward to apply per operator, since it asks each one to reason about the
block it belongs to. It becomes self-applying if the test is stated once, over the whole
block, and evaluated with the convention that an unbound anchor fails:
\[
\mathit{prune}_B(\mu) \;\equiv\; \bigwedge_{t \in \mathit{anchors}_V(B)}
\bigl( t \in \mathrm{dom}(\mu) \wedge \mu(t) \in CR \bigr).
\]
Register this one predicate --- the full residual anchor set of $B$, not the anchors of
the operator's own patterns --- on every operator covering a pattern of $B$. An operator
whose mappings do not bind some anchor of $B$ can never satisfy the conjunction and so
never prunes, which is the necessity direction of the lemma discharged automatically. An
operator that does bind them all prunes exactly the delegated part, which is the
sufficiency direction. No shape analysis is needed at registration time, and the same
predicate stays correct as the operator's coverage grows during adaptive execution.

The predicate must not fire on the mappings the resource itself supplies, which satisfy
every anchor condition by construction and would otherwise be deleted as though the
default plan had produced them. A mapping is exempt when any part of its derivation came
from the resource owning the block, which is a property of the derivation rather than of
its last step: in an adaptive plan the resource may equally have seeded a derivation or
been probed by one, and only an accumulated record of the resources a mapping passed
through distinguishes both cases.

\subsection*{Pruning Depth}

What the shape of $B$ decides is not \emph{whether} the test applies but \emph{how deep}
in the plan it first applies. Define the \emph{cover number} $\gamma(B)$ as the size of
the smallest $C \subseteq B$ with $\mathit{anchors}_V(B) \subseteq \mathit{vars}(C)$, and
the \emph{connected cover number} $\gamma_c(B)$ as the same quantity restricted to
connected $C$. An engine that does not form cartesian products materialises only
connected sub-blocks, so $\gamma_c$ is the operative one: by the Generalised Pruning
lemma it is exactly the least depth at which an operator of the default plan may prune.

\begin{proposition}
For a path of length $n$, $\gamma(B) = \lceil n/2 \rceil$ and $\gamma_c(B) = n - 1$. For
a star of $k$ patterns, $\gamma(B) = \gamma_c(B) = 1$.
\end{proposition}

\begin{proof}
Write $o_0 = {?s}$, so $\mathit{anchors}_V(B) = \{o_0, \dots, o_{n-1}\}$ and
$\mathit{vars}(tp_i) = \{o_{i-1}, o_i\}$. Each pattern contains at most two anchors,
giving $\gamma(B) \geq \lceil n/2 \rceil$, and $\{tp_1, tp_3, tp_5, \dots\}$ attains it.
A connected $C$ is a contiguous run $tp_i, \dots, tp_j$ with $\mathit{vars}(C) =
\{o_{i-1}, \dots, o_j\}$; covering $o_0$ forces $i = 1$ and covering $o_{n-1}$ forces
$j \geq n-1$, so $C \supseteq \{tp_1, \dots, tp_{n-1}\}$, which is itself a cover. For a
star the anchor set is a singleton contained in every pattern.
\end{proof}

\begin{corollary}[Work removed]
\label{cor:work}
Forming a smallest connected cover costs $\gamma_c(B) - 1$ joins, and evaluating $B$
costs $|B| - 1$, so pruning removes at most $|B| - \gamma_c(B)$ joins per delegated
mapping, attained when the plan forms a smallest connected cover. For a star this is
$k - 1$, the whole block. For a path it is $1$, independent of $n$.
\end{corollary}

The bound is not always attained, because an adaptive plan chooses its own join order. Of
the two connected sub-blocks of size $n-1$ in a path --- the prefix $tp_1 \dots tp_{n-1}$
and the suffix $tp_2 \dots tp_n$ --- only the prefix contains $?s$ and hence covers the
anchors. An order that builds the suffix first reaches the anchor cover only at the root,
where no join is left to skip and pruning coincides with deduplication.

\begin{center}
\begin{tabular}{llll}
\toprule
block & $\gamma_c$ & shallowest prunable operator & joins removed \\
\midrule
star at $?x$, $k$ patterns & $1$   & any single pattern               & $k - 1$ \\
path, $n = 2$              & $1$   & $tp_1$ alone                     & $1$ (all) \\
path, $n \geq 3$           & $n-1$ & the prefix $tp_1 \dots tp_{n-1}$ & $1$ of $n-1$ \\
\bottomrule
\end{tabular}
\end{center}

The two path cases differ in kind, not degree. At $n = 2$ the smallest cover has depth
$1$, so the test runs on $tp_1$ before its tuples are indexed: the delegated part never
enters the join state, and the reduction is a reduction of \emph{input}. At $n \geq 3$
the smallest connected cover has $n - 1$ patterns, so the test can only run after the
block has already been joined almost to completion: the delegated part is built and then
dropped, and the reduction is a reduction of \emph{output}. The delegated mappings still
consume hash-table space, still generate intermediate results, and still drive whatever
link discovery those results imply. What is saved is the final join and everything
downstream of it, a fraction $1/(n-1)$ of the block that vanishes as the path grows.

Applying the test at maximal depth, where the covering operator is $B$ itself, is
deduplication: pruning and deduplication are the same predicate evaluated at different
depths. This is also why the test cannot simply be omitted for long paths. If no
operator of $B$ ever evaluates it, the default plan reproduces the delegated part in
full, the four-way decomposition stops being a partition of what the engine emits, and
$CR$ contributes duplicates rather than results.

\section{The Scope of Delegation}

The pruning analysis gives a criterion for which blocks are worth delegating at all, and
it is sharper than a restriction on shape. We state it, show that it excludes long paths,
and show that long paths can nevertheless be delegated by decomposing them into blocks
that satisfy it.

\begin{definition}[Reducible block]
A block $B$ is \emph{reducible} if $\gamma_c(B) = 1$, that is, if some single pattern of
$B$ binds every residual anchor of $B$:
\[
\exists\, tp \in B \;:\; \mathit{anchors}_V(B) \subseteq \mathit{vars}(tp).
\]
\end{definition}

Reducibility is a property of the block, not of a shape vocabulary. It admits stars of
any size, whose anchor set is a singleton occurring in every pattern; paths of length
two, whose anchors $\{?s, ?o_1\}$ are the variables of $tp_1$; and any block, of whatever
shape, whose anchors happen to be collected in one pattern. It excludes paths of length
three or more, whose anchors are distributed across patterns.

\begin{theorem}[Characterisation]
The delegated part $B(A)$ contributes nothing to the join state of the default plan if
and only if $B$ is reducible.
\end{theorem}

\begin{proof}
By the Generalised Pruning lemma an operator may prune exactly when its mappings bind
every anchor of $B$, and by definition of $\gamma_c$ the shallowest such operator has
depth $\gamma_c(B)$. If $\gamma_c(B) = 1$ that operator scans a single pattern, so the
test is decided on each tuple as it is read, before it is indexed: no mapping of $B(A)$
is ever inserted into the join state. If $\gamma_c(B) > 1$ no depth-$1$ operator can
decide the test, so every tuple of $B(A)$ is indexed, and the mappings composing $B(A)$
are built by at least $\gamma_c(B) - 1$ joins before any of them can be discarded.
\end{proof}

The two cases differ in what pruning removes. At depth $1$ a discarded tuple never enters
the index, never probes, and never produces the intermediate results that would have
descended from it, so the saving is the whole sub-derivation rooted at that tuple and
scales with its fanout. At depth $\gamma_c(B) > 1$ the mapping has already been built,
and discarding it removes one join and the routing that would have followed. The first
saving is multiplicative in the data, the second is additive.

\subsection{Longer Paths Do Not Pay}

For a path, the excluded case can be quantified exactly. Write $D = B(A)$ for the
delegated part and $P$ for the mappings over a smallest connected cover that satisfy the
anchor conditions; each $p \in P$ extends to one or more members of $D$, so
$|P| \leq |D|$. By Corollary~\ref{cor:work} at most one join per delegated mapping is
removed, and only in the orders that form such a cover.

\begin{proposition}[Cost of delegating a path of length $n \geq 3$]
Let $c_{\mathit{req}}$ be the cost of a request, $c_{\mathit{parse}}$ the cost of
admitting one mapping from the resource, $c_{\mathit{join}}$ one join probe and
$c_{\mathit{mat}}$ one materialised mapping. Then
\[
W_{CR} - W_{\mathit{base}} \;=\;
c_{\mathit{req}} + |D|\,c_{\mathit{parse}}
\;-\; \bigl( |P|\,c_{\mathit{join}} + |D|\,c_{\mathit{mat}} \bigr),
\]
where the subtracted term is present only for the mappings whose join order forms a
smallest connected cover, and is zero otherwise.
\end{proposition}

\begin{proof}
Everything downstream of $B$ is unchanged: $|D|$ mappings enter the rest of the query
whether the resource or the default plan produced them. Within $B$, the default plan
performs the same joins as without the resource up to depth $\gamma_c(B) = n-1$, since no
shallower operator can prune. For the mappings that reach a cover, the final join and the
materialisation of their completions are skipped; for the others nothing is skipped. The
resource adds one request and the admission of $|D|$ mappings.
\end{proof}

Since $c_{\mathit{parse}} \geq c_{\mathit{mat}}$ --- admitting a mapping off the wire is
not cheaper than producing one from a hash probe --- the difference is non-negative in
every join order. Delegating a path of length three or more therefore never reduces total
work. What it reduces is latency: the default plan needs $n$ dependent round trips, since
the document binding $?o_{i}$ cannot be fetched before $?o_{i-1}$ resolves, whereas the
resource answers in one. The additional work is performed while the engine would
otherwise be blocked, so it is free in wall-clock terms exactly while the engine is
I/O-bound, and not otherwise.

\subsection{Reducible Decompositions}

Excluding long paths from delegation does not mean answering them without composite
resources, because a path can be partitioned into reducible blocks.

\begin{proposition}[Tiling]
Every connected sub-block of a path $tp_1, \dots, tp_n$ is a contiguous run, and a run of
$m$ patterns is reducible if and only if $m \leq 2$. Consequently every path admits a
partition into reducible blocks, and the smallest such partition has $\lceil n/2 \rceil$
blocks.
\end{proposition}

\begin{proof}
A run $tp_i, \dots, tp_j$ of length $m = j - i + 1$ has anchors
$\{o_{i-1}, \dots, o_{j-1}\}$ and $\gamma_c = m - 1$ for $m \geq 2$ by the earlier
proposition, and $\gamma_c = 1$ for $m = 1$; so it is reducible exactly when $m \leq 2$.
A partition into reducible blocks therefore uses blocks of at most two patterns and needs
at least $\lceil n/2 \rceil$ of them, and the partition into consecutive pairs, with one
singleton when $n$ is odd, attains that bound.
\end{proof}

The tiles are pattern-disjoint, so no two of them claim the same join entry, and each is
delegated under its own anchor conditions --- possibly to different pods, when the path
crosses pod boundaries. Their queries are parameterised rather than bound, so they do not
depend on one another and are requested concurrently: latency remains one round trip, as
it would be for a single resource covering the whole path, while every tile prunes at
depth $1$ and so removes its part of the delegated data from the default plan's join
state.

The tiles are, however, less selective than the path they decompose. A resource answering
$tp_1 \dots tp_4$ returns the mappings of the whole path, whereas two tiles return the
mappings of $tp_1 tp_2$ and of $tp_3 tp_4$ separately and leave their join to the client.
Decomposition therefore trades transferred volume and client-side join work for input
pruning and lower per-request server compute. Which side wins depends on the selectivity
of the path relative to its halves, and is an empirical question rather than one the
analysis settles.

\subsection{Guidance for Publishers}

Because reducibility is a property of the published block rather than of the client, the
criterion also says what a pod should expose. A resource answering a long path gives
clients a latency improvement and no reduction in work, since its result is rebuilt by
the client's default operators before being discarded. The same data published as
reducible tiles gives both, at the cost of one request per tile and a less selective
answer. A pod that publishes composite resources for its data is therefore better served
by publishing blocks of unit cover number --- stars over its subjects, and two-hop paths
--- than by publishing the longest paths it can compute.

The same argument applies to scope. A resource whose authority scope is a whole
namespace can be pruned against; one that aggregates a fragment of a namespace without
saying so cannot, and its clients must fall back to delegating without pruning. Declaring
the scope costs the publisher nothing and is what turns a resource from a latency
optimisation into a work reduction.

\subsection{What the Analysis Predicts}

The account above is falsifiable on three counts, and an evaluation that delegates long
paths as an ablation can check each directly:

\begin{enumerate}
\item time to first result improves for delegated paths of any length, in proportion to
      the number of dependent round trips removed;
\item total work increases for paths of length three or more, and decreases for stars and
      two-hop paths;
\item the distribution of pruning depth concentrates at depth $1$ for reducible blocks and
      at or near the root for longer paths, with the fraction reaching depth $n-1$
      governed by how often the adaptive order builds the prefix rather than the suffix.
\end{enumerate}

The third is the most informative, as it separates the two mechanisms: a long path whose
prunes all land at the root is contributing nothing but a request and a parse, and the
measurement says so without reference to the timings.


\end{document}
